\section{Proposed Method: Trust-Enhanced Graph-Based Recommendation (GTERM) System}

\subsection{Data Transformation through Trust-Augmented Graph Construction}

The transformation of interaction datasets into a graph representation is pivotal in GTERM, embedding trust metrics for enhanced recommendation quality within the Graph Neural Network (GNN) framework.

\textbf{Input:} Compiled from multiple datasets, including Yelp, Gowalla, ML-10M, and Amazon Book Ratings, providing a comprehensive user-preference and interaction history for graph construction.

\textbf{Process:}
\begin{enumerate}
    \item \textbf{Normalization and Cleanup:} Data normalization (min-max scaling, Z-score) ensures uniformity in interaction scores across various datasets.
    
    \item \textbf{Trust Metric Derivation:} Trust scores calculated from historical interaction metrics (e.g., review frequencies, ratings) are encoded as edge weights. This enhances the structural scope of the graph with trust measures bolstering the framework's interpretative strength.
    
    \item \textbf{Graph and Node Configuration:} The bipartite graph \(G = (U \cup I, E)\), with \(U\) and \(I\) as user and item sets respectively, utilizes trust-inflected edges \(E\). Node attributes are initialized with vectors capturing both interaction history and trust levels.

    \item \textbf{Efficient Processing Adaptation:} PyTorch Geometric DataLoader partitions graph data into manageable subsets, optimizing memory and computational efficiency for GNN workflows.

    \item \textbf{Robustness Mechanism:} Detailed logging and diagnostic checks are implemented to safeguard against data handling discrepancies, ensuring robustness during graph convolutions.
\end{enumerate}


\textbf{Implementation:} Implemented in the module \texttt{process\_data.py}, leveraging PyTorch Geometric for sparse matrix manipulation, incorporating trust-enhanced connectivity. Future iterations will include attention mechanisms to prioritize higher-trust interactions for refined recommendation precision.


\textbf{Output:} The output is a dynamic graph construct ready for GNN integration, promoting accuracy through trust-aware interaction mappings.

\subsection{GNN-Based Interaction Prediction}

GNNs in GTERM facilitate intricate user-item interaction modeling, with trust metrics enhancing the accuracy and fidelity of these predictions.

\paragraph{Architecture} A hybrid approach utilizing Graph Convolutional Networks (GCN) and Graph Attention Networks (GAT) refines node embeddings:

\begin{itemize}
    \item \textbf{GCN Layer:} Aggregates neighborhood information via:
    \[
    \mathbf{H}^{(1)} = \sigma(\widehat{\mathbf{D}}^{-\frac{1}{2}} \widehat{\mathbf{A}} \widehat{\mathbf{D}}^{-\frac{1}{2}} \mathbf{X} \mathbf{W}^{(0)})
    \]
    ensuring robust feature propagation and structural dependency modeling, as seen in \cite{kipf2017semisupervised}.

    \item \textbf{GAT Layer:} Employs attention to emphasize trust-influenced interactions:
    \[
    \alpha_{ij} = \frac{\exp(\text{LeakyReLU}(\mathbf{a}^\top[\mathbf{W}^{(1)} \mathbf{H}^{(1)}_i \, \| \, \mathbf{W}^{(1)} \mathbf{H}^{(1)}_j]))}{\sum_{k \in \mathcal{N}(i)} \exp(\text{LeakyReLU}(\mathbf{a}^\top[\mathbf{W}^{(1)} \mathbf{H}^{(1)}_i \, \| \, \mathbf{W}^{(1)} \mathbf{H}^{(1)}_k]))}
    \]
    advancing the model's selectivity on significant graph edges \cite{velivckovic2018graph}.
\end{itemize}

\paragraph{Integration} A global pooling operation refines node embeddings into graph representations, processed through dense layers to yield interaction probabilities.

\paragraph{Composite Score Synthesis} Combining interaction predictions with trust:
\[
\text{Composite Score} = \alpha \times \text{Trust Score} + (1 - \alpha) \times \text{GNN Score}
\]
\(\alpha\) determines the weight between interaction and trust, optimizing recommendation transparency and robustness.

\subsubsection{Trust-Weighted Interaction Prioritization}

Enhancing interaction scores by weighting trust, aligning with frameworks seen in similar domains.

\subsection{Transparent Trust-Integrated Recommendation Framework}

Integrating trust metrics within GTERM augments both recommendation precision and interpretability.

\textbf{Method Overview:}

\begin{enumerate}
    \item \textbf{Graph Construction:} Trust-score-weighted edges provide a semantic information layer within the bipartite graph.
    
    \item \textbf{Weighted Interaction Scoring:} Embeddings honed through GNN layers reflect trust dynamics and adapt according to significant interactions.

    \item \textbf{Prediction and Trust Integration:} Composite scores meld GNN predictions with trust inputs, balancing accuracy with user transparency.

    \item \textbf{Explanatory Path Output:} Every recommendation accompanies a path-based explanation derived from graph structures, elucidating the decision-making rationale.
\end{enumerate}

\textbf{Benefits} This method amplifies recommendation credibility, driven by a seamless integration of advanced GNN insights and comprehensive trust factors, enhancing effectiveness and user confidence.
